% Options for packages loaded elsewhere
\PassOptionsToPackage{unicode}{hyperref}
\PassOptionsToPackage{hyphens}{url}
\PassOptionsToPackage{table}{xcolor}
%
\documentclass[
  11pt,
]{report}
\usepackage{amsmath,amssymb}
\usepackage{iftex}
\ifPDFTeX
  \usepackage[T1]{fontenc}
  \usepackage[utf8]{inputenc}
  \usepackage{textcomp} % provide euro and other symbols
\else % if luatex or xetex
  % pdfLaTeX-only repository policy: the original unicode-math/fontspec
  % branch is unreachable under pdfLaTeX and is replaced with a guard.
  \PackageError{cpython-suite}{Compile with pdfLaTeX only}{}
\fi
\usepackage{lmodern}
\ifPDFTeX\else
  % xetex/luatex font selection
\fi
% Use upquote if available, for straight quotes in verbatim environments
\IfFileExists{upquote.sty}{\usepackage{upquote}}{}
\IfFileExists{microtype.sty}{% use microtype if available
  \usepackage[]{microtype}
  \UseMicrotypeSet[protrusion]{basicmath} % disable protrusion for tt fonts
}{}
\makeatletter
\@ifundefined{KOMAClassName}{% if non-KOMA class
  \IfFileExists{parskip.sty}{%
    \usepackage{parskip}
  }{% else
    \setlength{\parindent}{0pt}
    \setlength{\parskip}{6pt plus 2pt minus 1pt}}
}{% if KOMA class
  \KOMAoptions{parskip=half}}
\makeatother
\usepackage{xcolor}
\usepackage{listings}
\newcommand{\passthrough}[1]{#1}
\lstset{defaultdialect=[5.3]Lua}
\lstset{defaultdialect=[x86masm]Assembler}
\usepackage{longtable,booktabs,array}
\usepackage{calc} % for calculating minipage widths
% Correct order of tables after \paragraph or \subparagraph
\usepackage{etoolbox}
\makeatletter
\patchcmd\longtable{\par}{\if@noskipsec\mbox{}\fi\par}{}{}
\makeatother
% Allow footnotes in longtable head/foot
\IfFileExists{footnotehyper.sty}{\usepackage{footnotehyper}}{\usepackage{footnote}}
\makesavenoteenv{longtable}
\setlength{\emergencystretch}{3em} % prevent overfull lines
\providecommand{\tightlist}{%
  \setlength{\itemsep}{0pt}\setlength{\parskip}{0pt}}
\setcounter{secnumdepth}{5}
% CPython Engineering Suite shared presentation layer -- pdfLaTeX compatible
\usepackage[T1]{fontenc}
\usepackage{lmodern}
\usepackage{microtype}
\usepackage{geometry}
\geometry{letterpaper,margin=0.91in,headheight=16pt,headsep=19pt}
\usepackage{xcolor} % [table] passed at top of file (option-clash fix)
\definecolor{CPNavy}{HTML}{13283B}
\definecolor{CPBlue}{HTML}{176BA0}
\definecolor{CPGold}{HTML}{D8A73C}
\definecolor{CPMuted}{HTML}{566979}
\definecolor{CPBg}{HTML}{EFF5F8}
\usepackage{titlesec}
\titleformat{\chapter}[display]{\normalfont\sffamily\bfseries\color{CPNavy}}{\large\color{CPBlue}CHAPTER \thechapter}{9pt}{\Huge}[\vspace{6pt}\color{CPGold}\titlerule]
\titleformat{\section}{\Large\bfseries\sffamily\color{CPBlue}}{\thesection}{0.7em}{}
\titleformat{\subsection}{\large\bfseries\sffamily\color{CPNavy}}{\thesubsection}{0.7em}{}
\usepackage{fancyhdr}
\pagestyle{fancy}
\fancyhf{}
\fancyhead[L]{\footnotesize\color{CPMuted}\sffamily CPYTHON ENGINEERING SUITE}
\fancyhead[R]{\footnotesize\color{CPMuted}\sffamily 3.16.0a1+dev SNAPSHOT}
\fancyfoot[C]{\footnotesize\color{CPMuted}\thepage}
\renewcommand{\headrulewidth}{0.3pt}
\AtBeginDocument{\urlstyle{tt}}
\lstset{basicstyle=\ttfamily\small,columns=fullflexible,keepspaces=true,breaklines=true,breakatwhitespace=false,showstringspaces=false,frame=single,framesep=5pt,rulecolor=\color{CPBg},backgroundcolor=\color{CPBg},xleftmargin=3pt,xrightmargin=3pt}
\usepackage{tcolorbox}
\tcbuselibrary{breakable,skins}
\newtcolorbox{sourcebox}[1][]{enhanced,breakable,sharp corners,boxrule=0.4pt,colback=CPBg,colframe=CPBlue,borderline west={3pt}{0pt}{CPBlue},left=10pt,right=8pt,top=6pt,bottom=6pt,#1}
\newtcolorbox{cautionbox}[1][]{enhanced,breakable,sharp corners,boxrule=0.4pt,colback=CPBg,colframe=CPGold,borderline west={3pt}{0pt}{CPGold},left=10pt,right=8pt,top=6pt,bottom=6pt,#1}
\usepackage{tikz}
\usetikzlibrary{arrows.meta,positioning,shapes.geometric,calc}
\usepackage{enumitem}
\setlist{itemsep=2pt,topsep=4pt}
\usepackage{booktabs}
\usepackage{longtable}
\AtBeginEnvironment{longtable}{\small}
\usepackage{array}
\setlength{\parskip}{5pt}
\setlength{\parindent}{0pt}
\setcounter{tocdepth}{2}
\usepackage{hyperref}
\hypersetup{colorlinks=true,linkcolor=CPBlue,urlcolor=CPBlue,citecolor=CPBlue,pdfauthor={CPython Engineering Documentation Suite}}
\makeatletter
\renewcommand{\maketitle}{%
  \begin{titlepage}\pagecolor{CPNavy}\color{white}\sffamily
  \vspace*{0.8in}
  {\color{CPGold}\rule{2.1in}{4pt}}\\[16pt]
  {\Large\bfseries CPYTHON ENGINEERING DOCUMENTATION SUITE}\\[24pt]
  {\Huge\bfseries \@title\par}\vspace{19pt}
  \vfill
  {\normalsize Source baseline: CPython 3.16.0 alpha 1+dev repository snapshot}\par
  {\normalsize Prepared 16 October 2026}\par
  {\small\color{white!70} Source-traced engineering commentary; not an official Python specification.}\vspace{0.65in}
  \end{titlepage}\nopagecolor\color{black}
  }
\makeatother
\ifLuaTeX
  \usepackage{selnolig}  % disable illegal ligatures
\fi
\usepackage{bookmark}
\IfFileExists{xurl.sty}{\usepackage{xurl}}{} % add URL line breaks if available
\urlstyle{same}
\hypersetup{
  pdftitle={Python C API and Extension Engineering},
  hidelinks,
  pdfcreator={LaTeX via pandoc}}

\title{Python C API and Extension Engineering}
\usepackage{etoolbox}
\makeatletter
\providecommand{\subtitle}[1]{% add subtitle to \maketitle
  \apptocmd{\@title}{\par {\large #1 \par}}{}{}
}
\makeatother
\subtitle{Volume 06 \textbar{} Ownership, modules, embedding, ABI
stability and native safety}
\author{}
\date{}

\begin{document}
\maketitle

{
\setcounter{tocdepth}{1}
\tableofcontents
}
\chapter{Public C API contract and stability
tiers}\label{public-c-api-contract-and-stability-tiers}

\begin{sourcebox}
\textbf{Evidence classification.} \textbf{SOURCE} = statements supported by material present in the uploaded snapshot; \textbf{IMPLEMENTATION} = CPython-specific internals, not guaranteed Python semantics; \textbf{ENGINEERING} = original recommendations, examples, checks, and exercises. The source baseline is \textbf{3.16.0 alpha 1+dev}; missing repository file bodies are noted as evidence limitations.
\end{sourcebox}

CPython provides native integration for extension modules and embedding.
\path{Doc/c-api/intro.rst} describes these as the two
broad ways to use the Python/C API.
\path{Doc/c-api/stable.rst} discusses backward
compatibility and ABI expectations, including the \textbf{Limited API},
\textbf{Stable ABI} and \textbf{Unstable API} (the latter with
\passthrough{\lstinline!PyUnstable!} prefixes). Private implementation
facilities, particularly \passthrough{\lstinline!\_Py!} functions and
\path{Include/internal/}, must not be mistaken for
stable extension contracts.

\textbf{Version:} the uploaded repository is \textbf{3.16.0 alpha
1+dev}. \path{Doc/c-api/extension-modules.rst}
contains a \path{versionadded:: 3.15} note for the
\path{PyModExport_*} export hook while documenting
older-style \passthrough{\lstinline!PyInit\_*!} entry points as
available. Select the export interface only after verifying the minimum
supported interpreter version and module tooling.

\begin{longtable}[]{@{}
  >{\raggedright\arraybackslash}p{(\columnwidth - 4\tabcolsep) * \real{0.3333}}
  >{\raggedright\arraybackslash}p{(\columnwidth - 4\tabcolsep) * \real{0.3333}}
  >{\raggedright\arraybackslash}p{(\columnwidth - 4\tabcolsep) * \real{0.3333}}@{}}
\toprule\noalign{}
\begin{minipage}[b]{\linewidth}\raggedright
Native surface
\end{minipage} & \begin{minipage}[b]{\linewidth}\raggedright
Use case
\end{minipage} & \begin{minipage}[b]{\linewidth}\raggedright
Change exposure
\end{minipage} \\
\midrule\noalign{}
\endhead
\bottomrule\noalign{}
\endlastfoot
Public C API & Extensions targeting a specific supported CPython version
& Subject to documented compatibility policy \\
Limited API & Source code restricted to an API subset & Intended for
wider source version reach \\
Stable ABI / \passthrough{\lstinline!abi3!} & Binary compatible
extension distribution within documented limits & Subject to
platform/build constraints and approved symbol set \\
\passthrough{\lstinline!PyUnstable\_*!} & Special runtime integrations
with explicit instability & May change more frequently \\
Private \passthrough{\lstinline!\_Py*!} / internal headers & CPython
development & No ordinary third-party stability promise \\
\end{longtable}

The Stable ABI does not automatically make arbitrary extensions
subinterpreter-safe, thread-safe or independent of native library ABI
dependencies.

\chapter{Object ownership and error
contracts}\label{object-ownership-and-error-contracts}

A \passthrough{\lstinline!PyObject*!} normally denotes a reference to a
Python object; the \textbf{reference ownership contract} is critical.
Strong (owned) references must be decremented by the owner, while
borrowed references require another owner to remain valid. An API that
steals a reference has a different transfer contract. Always consult the
precise function's documentation instead of inferring ownership from its
name alone.

\begin{lstlisting}[language=C]
#include <Python.h>

/* Create a tuple containing a strong reference to argument. */
static PyObject *make_pair(PyObject *value)
{
    PyObject *pair = PyTuple_New(2);
    if (pair == NULL) {
        return NULL;
    }
    PyObject *left = Py_NewRef(value);
    PyObject *right = Py_NewRef(value);
    /* PyTuple_SET_ITEM steals each new reference. */
    PyTuple_SET_ITEM(pair, 0, left);
    PyTuple_SET_ITEM(pair, 1, right);
    return pair;
}
\end{lstlisting}

The example is for API ownership training, not a recommendation to
bypass all bounds checks in general-purpose extension code.
\path{PyTuple_SET_ITEM} and its preconditions must
be reviewed against the supported API tier. The
\passthrough{\lstinline!NULL!} return convention indicates Python error
propagation for many \passthrough{\lstinline!PyObject*!} returning C API
functions. Never return \passthrough{\lstinline!NULL!} without an active
exception unless the specific API contract expects it.

\section{Defensive error-path
pattern}\label{defensive-error-path-pattern}

\begin{lstlisting}[language=C]
static PyObject *checked_double(PyObject *self, PyObject *arg)
{
    (void)self;
    PyObject *two = PyLong_FromLong(2);
    if (two == NULL) return NULL;
    PyObject *result = PyNumber_Multiply(arg, two);
    Py_DECREF(two);
    return result; /* on failure, propagates the existing exception */
}
\end{lstlisting}

\textbf{ENGINEERING review:} for each allocation, record the owner; for
each operation that can fail, check its return value; for each exit,
inspect whether all owned references have been released or intentionally
transferred.

\chapter{A complete introductory extension
module}\label{a-complete-introductory-extension-module}

The following is an \textbf{educational} single-phase initialization
example using widely understood APIs. It demonstrates a native method,
\passthrough{\lstinline!METH\_O!}, error propagation and module
registration. It is not, by itself, a certification of free-threaded or
multi-interpreter compatibility.

\begin{lstlisting}[language=C]
#define PY_SSIZE_T_CLEAN
#include <Python.h>

static PyObject *twice(PyObject *self, PyObject *value)
{
    (void)self;
    PyObject *two = PyLong_FromLong(2);
    if (two == NULL) return NULL;
    PyObject *result = PyNumber_Multiply(value, two);
    Py_DECREF(two);
    return result;
}

static PyMethodDef methods[] = {
    {"twice", (PyCFunction)twice, METH_O, "Return value * 2."},
    {NULL, NULL, 0, NULL}
};
static struct PyModuleDef module = {
    .m_base = PyModuleDef_HEAD_INIT,
    .m_name = "demo_native",
    .m_doc = NULL,
    .m_size = -1,
    .m_methods = methods
};
PyMODINIT_FUNC PyInit_demo_native(void)
{
    return PyModule_Create(&module);
}
\end{lstlisting}

\textbf{ENGINEERING build guidance:} use a maintained packaging build
backend for distributions rather than pasting an unreviewed compiler
command into production build automation. Verify that the compiled
module's suffix, ABI, toolchain and architecture match the chosen Python
installation. For this example, import the built module and test
\path{demo_native.twice(21) == 42} plus a deliberate
unsupported-operand failure.

\chapter{Multi-phase module initialization and state
isolation}\label{multi-phase-module-initialization-and-state-isolation}

\path{Doc/c-api/extension-modules.rst} and
\path{Doc/c-api/module.rst} discuss module
initialization and state. For modules intended to support multiple
interpreters and free-threaded builds, put mutable state in
module-associated storage rather than unchecked process-wide globals. A
state design must address traversal, clearing, freeing and any native
resources held beyond module lifetime.

\textbf{Implementation decision record:}

\begin{enumerate}
\def\labelenumi{\arabic{enumi}.}
\tightlist
\item
  Is compatibility required with Python versions older than 3.15?
\item
  Is a Stable ABI / \passthrough{\lstinline!abi3!} wheel required, and
  which Limited API subset is permitted?
\item
  Will modules run in multiple interpreters in the same process?
\item
  Is the free-threaded build supported? If so, where is synchronization
  enforced?
\item
  Does the module retain pointers into movable/reallocatable buffers or
  borrowed references?
\item
  Can a module be initialized, used, cleared and torn down more than
  once?
\end{enumerate}

The snapshot documents a
\path{PyModExport_modulename} hook added in 3.15;
the exact slots and declarations must follow the corresponding version's
\path{Doc/c-api/extension-modules.rst}, not be
inferred from a legacy \passthrough{\lstinline!PyInit\_*!} example.

\chapter{The buffer protocol, native memory and
callbacks}\label{the-buffer-protocol-native-memory-and-callbacks}

The buffer protocol enables borrowing or exporting access to memory,
with exporter/importer lifetime contracts. Mismanagement can cause
use-after-free, stale pointer access, races or crashes. Releasing the
GIL around native work may improve concurrency for I/O or expensive
native operations, but native code must not touch Python objects
unsafely while detached from the relevant thread state.

\textbf{Security review checklist:} validate lengths before arithmetic,
use \passthrough{\lstinline!Py\_ssize\_t!} for Python-size quantities,
check signed/unsigned conversion, defend against integer overflow before
allocation, avoid returning pointers to stack-local memory, and ensure
all temporary buffers are released on exception paths. For external
buffers, document whether the buffer may be mutated during a native
call.

\chapter{Embedding Python in a host
process}\label{embedding-python-in-a-host-process}

\path{Doc/c-api/init.rst},
\path{Doc/c-api/init_config.rst} or relevant
initializer documentation, and
\path{Doc/extending/embedding.rst} describe
integration from the host's perspective. Treat Python startup as
configuration management: module search paths, environment handling,
process state, locale, logging, audit and filesystem authority are
operational security considerations. A host that executes arbitrary
embedded Python code is not sandboxed by using the C API.

\textbf{ENGINEERING execution sequence:}

\begin{lstlisting}
host policy + runtime config
          -> initialize Python
          -> configure import locations / modules
          -> acquire correct interpreter and thread state
          -> call Python API with checked exception results
          -> release references and native resources
          -> finalize in defined host shutdown order
\end{lstlisting}

Do not assume finalization may safely occur while unrelated native
threads still use the runtime. Embedding uses a larger lifecycle than an
extension method called by an already-running interpreter.

\chapter{Free-threading and subinterpreter
readiness}\label{free-threading-and-subinterpreter-readiness}

\path{Doc/howto/free-threading-extensions.rst}
identifies the \path{Py_GIL_DISABLED} macro in
free-threaded builds and explains how extensions declare compatibility.
A module that omits the declaration may cause the runtime to enable the
GIL. Declaration should follow an actual concurrency audit. Modules with
process-global caches, borrowed-reference sharing or unsynchronized
state require special attention.
\path{Doc/c-api/subinterpreters.rst} covers
interpreter separation requirements.

\textbf{Testing matrix:}

\begin{longtable}[]{@{}
  >{\raggedright\arraybackslash}p{(\columnwidth - 2\tabcolsep) * \real{0.5000}}
  >{\raggedright\arraybackslash}p{(\columnwidth - 2\tabcolsep) * \real{0.5000}}@{}}
\toprule\noalign{}
\begin{minipage}[b]{\linewidth}\raggedright
Configuration
\end{minipage} & \begin{minipage}[b]{\linewidth}\raggedright
Primary validation
\end{minipage} \\
\midrule\noalign{}
\endhead
\bottomrule\noalign{}
\endlastfoot
Normal CPython build & Correct API results, exceptions, ownership \\
Debug build & Assert/reference-sensitive behavior \\
Free-threaded build & Concurrent access, locks, shared-state
invariants \\
Multiple interpreters & No unauthorized cross-interpreter object
sharing \\
Minimum supported release & Availability of symbols and module export
scheme \\
Stable ABI build (if claimed) & Wheel tagging, symbol restriction and
import compatibility \\
\end{longtable}

\chapter{Hands-on laboratories}\label{hands-on-laboratories}

\section{Lab 1: Ownership and fault
injection}\label{lab-1-ownership-and-fault-injection}

Implement a new-reference-returning helper. Intentionally make one
allocation fail in a test harness; prove that all preceding owned
references are cleaned up and a Python exception is propagated.

\section{Lab 2: Extension function
tests}\label{lab-2-extension-function-tests}

Build \passthrough{\lstinline!demo\_native!}, test arithmetic inputs and
unsupported operands, then inspect native debug stack traces for a
failing call. Explain why the single-phase example is insufficient as
proof of multi-interpreter safety.

\section{Lab 3: ABI decision}\label{lab-3-abi-decision}

Write an engineering decision record comparing a version-specific
CPython extension, Limited API (\passthrough{\lstinline!abi3!})
extension and pure Python implementation. Include deployment versions,
required APIs, performance measurement plan and maintenance cost.

\section{Lab 4: Free-threading audit}\label{lab-4-free-threading-audit}

Inventory all mutable globals, reference ownership transfers, borrowed
pointers, allocation lifetime assumptions and thread entry/exit points
in a small extension. Produce a change list and reproducible thread
stress tests.

\chapter{Source concordance}\label{source-concordance}

\begin{itemize}
\tightlist
\item
  \textbf{SOURCE:} \path{Doc/c-api/intro.rst},
  \passthrough{\lstinline!stable.rst!},
  \path{refcounting.rst},
  \path{extension-modules.rst},
  \passthrough{\lstinline!module.rst!},
  \passthrough{\lstinline!threads.rst!},
  \passthrough{\lstinline!buffer.rst!},
  \passthrough{\lstinline!memory.rst!}.
\item
  \textbf{SOURCE:}
  \path{Doc/extending/first-extension-module.rst},
  \path{extending.rst},
  \path{embedding.rst};
  \path{Doc/howto/free-threading-extensions.rst},
  \path{abi3t-migration.rst}.
\item
  \textbf{SOURCE:} \path{Include/Python.h},
  \path{Include/object.h},
  \path{Include/refcount.h},
  \path{Include/pybuffer.h},
  \path{Include/moduleobject.h}.
\item
  \textbf{ENGINEERING:} examples, ownership ledger, architecture
  decisions and exercises. Neither the sample module nor the written
  checklist supersedes the deployed version's native API documentation.
\end{itemize}

\end{document}
